% HIKET -- storyline v2. A SKELETON, deliberately.
% Successor to HIKET_storyline_note.tex (retired 2026-09-02).
%
% PURPOSE. Lay the story out as headers first, so the ARC can be argued about
% before any of it is written. Each block carries one to three sentences and the
% figure it rests on -- nothing more. Prose belongs in the manuscript, not here.
%
% WHAT IS NEW vs THE OLD NOTE. Two things:
%   (1) Thread A is restated. "Complexity does not help" has become "fit does not
%       discriminate structure, but dynamics do" -- a positive claim, and the one
%       that matters for applied use.
%   (2) A new central beat: the calibration wants MORE INPUT and FASTER KINETICS
%       than the published parameterisations, and that joint displacement should
%       have consequences we can TEST -- starting with climate reactivity.
\documentclass[11pt]{article}
\usepackage[utf8]{inputenc}\usepackage[T1]{fontenc}
\usepackage[margin=2.4cm]{geometry}
\usepackage{graphicx}\usepackage{amsmath}\usepackage{amssymb}\usepackage{xcolor}
\usepackage{float}\usepackage[hidelinks]{hyperref}
\usepackage{caption}\captionsetup{font=small,labelfont=bf}
\usepackage[most]{tcolorbox}
\usepackage{titlesec}
\titlespacing*{\paragraph}{0pt}{0.9\baselineskip}{0.4em}
\setcounter{secnumdepth}{4}\setcounter{tocdepth}{3}

\definecolor{notebg}{HTML}{EAF2FB}\definecolor{noterule}{HTML}{2C6FB5}
\newtcolorbox{note}{colback=notebg, colframe=noterule, boxrule=0pt, leftrule=3pt,
  arc=1pt, left=8pt, right=8pt, top=5pt, bottom=5pt, fontupper=\small,
  before skip=7pt, after skip=7pt}
\definecolor{testbg}{HTML}{EAF7EE}\definecolor{testrule}{HTML}{2E8B57}
\newtcolorbox{test}{colback=testbg, colframe=testrule, boxrule=0pt, leftrule=3pt,
  arc=1pt, left=8pt, right=8pt, top=5pt, bottom=5pt, fontupper=\small,
  before skip=7pt, after skip=7pt}
\definecolor{warnbg}{HTML}{FDF3E3}\definecolor{warnrule}{HTML}{C08A2E}
\newtcolorbox{warn}{colback=warnbg, colframe=warnrule, boxrule=0pt, leftrule=3pt,
  arc=1pt, left=8pt, right=8pt, top=5pt, bottom=5pt, fontupper=\small,
  before skip=7pt, after skip=7pt}

\newcommand{\sinit}{\sigma_{\mathrm{init}}}
\newcommand{\sinput}{\sigma_{\mathrm{input}}}
\newcommand{\FIG}[1]{\textcolor{noterule}{\small$\blacktriangleright$ \textsf{#1}}}
\newcommand{\NEW}[1]{\textcolor{testrule}{\small$\bigstar$ \textsf{TO BUILD: #1}}}

\title{\vspace{-1.4cm}\textbf{HIKET --- the storyline, v2}\\[2pt]
\large skeleton for discussion: headers, one-line arguments, figures}
\author{}
\date{\small Version of \today. Numbers keyed to RUN\_IDs \texttt{20260831\_1624*} (arm B).
Deliberately not prose --- the point is to argue about the ARC.}

\begin{document}\maketitle

\begin{note}
\textbf{How to read this.} Every block is a claim we intend to make, its evidence, and nothing
else. \FIG{name} marks a figure that exists; \NEW{name} marks one that does not yet.
Green boxes are \emph{tests we could run}; amber boxes are \emph{things that could sink the claim}.
\end{note}

\tableofcontents
\bigskip

% =====================================================================
\section{The arc in one paragraph}
Finnish forest soils are accumulating carbon and the accumulation is beginning to level off. A model
started at steady state cannot show that, because a steady state is already the end of it. Started
instead from a calibrated past, all six models reach the observed levels and two of them reproduce
the rate exactly --- but doing so requires \emph{more litter input and faster turnover} than the
published parameterisations carry --- in two of the three Yassos. That displacement is the finding,
and it is not free: it changes how the soil responds to warming, but \textbf{only where a single
climate modifier ties turnover and temperature sensitivity to the same parameter}. Where a model
carries pool-specific modifiers, the same displacement costs nothing, and the projection is
\emph{resilient} to an input assumption the data cannot settle. \textbf{That contrast, measured, is
the contribution.}

% =====================================================================
\part{Act I --- Opening}

\section{A sink caught saturating}
\paragraph{The observed dynamic.} On the balanced plot set the mean stock runs $64.5 \to 71.5 \to
73.6$~tC\,ha$^{-1}$ across the three campaigns: $+7.0$ then $+2.1$. Accumulation, bending over.
\FIG{F3\_mean\_soc}

\paragraph{The rate is what an inventory reports.} $+0.259 \pm 0.040$~tC\,ha$^{-1}$\,yr$^{-1}$ over
the full span. \FIG{F5\_stock\_change}

\paragraph{Why it is politically live.} The Finnish GHG inventory turns on whether this sink
persists or saturates. \FIG{F10b\_forest\_history}

\section{The convention that cannot see it}
\paragraph{Steady-state initialisation is near-universal and rarely examined.} Lehtonen 2016,
Palosuo 2008, Peltoniemi 2004 --- our own group's lineage, to be framed as evolving, not attacked.

\paragraph{An equilibrium start is already at the end.} It cannot produce a rise, so it cannot
produce a rise that is levelling off.

% =====================================================================
\part{Act II --- Challenge}

\section{To show the saturation, the model must remember it was rising}
\paragraph{The soil carries a memory the model is told to forget.} Slow pools integrate a century
of inputs; setting them to equilibrium erases exactly the transient the data show.

\paragraph{Managed landscapes are rarely at equilibrium.} Finland is the vivid case, not a special
one: exploitation, then a century of managed regrowth. \FIG{F10b\_forest\_history}

% =====================================================================
\part{Act III --- Action}

\section{Transient initialisation under a fair frame}
\paragraph{The unknown start becomes a calibrated parameter.} Each draw spins up from 1917;
$\sinit$ is the ratio of 1917 to 1985 litter input. Uncertainty propagates rather than being patched.
\FIG{F4\_initialization}

\paragraph{Fair means equal \emph{information}, not equal parameter count.} Yasso does not fit its
decomposition rates, so neither do the simple models: rates externally anchored (ICBM), fractions
free. Otherwise structure and calibration asymmetry are confounded.

\paragraph{One likelihood, and it does not pretend the plots are independent.} Log-normal with a
fixed correlated error structure (region, plot, campaign), total $\sigma = 0.800$ partitioned, not
added. \NEW{schematic of the variance partition --- optional}

% =====================================================================
\part{Act IV --- Resolution, part 1: it works}

\section{The models reach the observations}
\paragraph{Levels.} No model sits more than $3.5$~tC\,ha$^{-1}$ from any campaign mean. The old
``$26$~tC\,ha$^{-1}$ too high at 1985'' caveat is gone. \FIG{F3\_mean\_soc}

\paragraph{The rate.} Against observed $+0.259$: Yasso07 $+0.258$, Yasso15 $+0.250$. TP2/TP3 about a
third strong; SP1/Yasso20 about half weak. The ensemble brackets the answer.
\FIG{F5\_stock\_change}

\paragraph{\textbf{The three-way split --- new, and better than ``all six work''.}}
Yasso07/Yasso15 rise then flatten (\emph{saturation reproduced}); TP2/TP3 rise and keep rising
(\emph{no saturation}); SP1/Yasso20 overshoot then decline. Structure is legible in the shape.
\NEW{F3b --- trajectory shape, three groups, explicitly grouped}

\begin{warn}
The 2006--2024 interval discriminates far less than it looks: observed $+0.117 \pm 0.066$, CI
includes zero. Do not lean on ``two models get the sign wrong'' over eighteen years.
\end{warn}

% =====================================================================
\part{Act V --- Thread A, restated}

\section{Fit does not discriminate structure; dynamics do}
\begin{note}
\textbf{This replaces ``does complexity improve the answer?''.} The old answer (``no, and the most
complex fits worst'') is no longer true --- Yasso15 has the best calibration $R^2$ \emph{and} the
best holdout RMSE. The defensible and more useful claim is the one below.
\end{note}

\subsection{In-sample, structure is invisible}
\paragraph{Skill does not climb the pool ladder.} Calibration $R^2$ $0.016$--$0.021$; holdout $R^2$
$\le 0.0012$ in all six; holdout RMSE spans $3.6\%$. One pool and five pools are indistinguishable.
\FIG{T1\_model\_summary} \FIG{F6b\_skill\_vs\_complexity}

\paragraph{The ordering does not follow structure at all.} Best and worst holdout RMSE belong to a
five-pool and a three-pool model.

\subsection{Out of sample, structure is nearly everything}
\paragraph{The forecast diverges where the calibration agreed.} \FIG{F4\_initialization} (panel b)

\paragraph{Climate sensitivity differs $\sim$2$\times$ across the family.} Equilibrium SOC change at
$+2$\,\textdegree C: Yasso07 $-21.7\%$, Yasso20 $-13.7\%$, Yasso15 $-9.6\%$. Between-model range is
$3.1\times$ the mean within-model band. \FIG{climate\_sensitivity\_sweep}
\begin{warn}
Measured on an earlier run and as a \emph{demonstration} (step change, single-exponential on bulk
MRT). Must be recomputed on the current calibration before it goes in a figure.
\end{warn}

\subsection{The applied consequence}
\paragraph{Model choice is a forecast decision, not a fit decision.} Selecting an inventory model on
calibration skill selects on a quantity that carries almost no information about what the model will
project. This is the practical message of the paper.

\paragraph{So the deliverable is an ensemble, not a winner.} An estimate plus its structural
uncertainty.

% =====================================================================
\part{Act VI --- Thread B, the new centre}

\section{The calibration wants more input \emph{and} faster kinetics}
\begin{note}
\textbf{This is the turn, and it should land as a surprise.} The paper sets out to fix an
initialisation convention. The expected ending is the tidy one: the convention was the problem, we
replaced it, the models work. That ending is \emph{half} true --- and the other half is the
interesting part. Making the models reproduce the observed trajectory requires a soil that is
\textbf{fed more and turns over faster} than any of the published parameterisations describe. A
methods fix walks into a claim about the models themselves.
\end{note}

\subsection{Why the surprise is the method's own doing}
\paragraph{A steady-state start hides exactly this.} At equilibrium $C = J/k$: only the
\emph{product} matters, so any input-and-turnover pair with the right product reproduces the stock.
Nothing in the fit can object to the published pair.

\paragraph{A transient start adds timing, and timing separates them.} Once the model must
\emph{travel} from 1917 through three observed campaigns, the rate of approach matters, and that is
set by the turnover time on its own. The constraint the convention discarded is the one that moves
the ensemble off the published point.

\paragraph{So the twist is not incidental to the method --- it is caused by it.} The initialisation
fix is what makes the parameter displacement visible. That is why the two halves belong in one paper.

\begin{warn}
\textbf{Do not overstate the separation.} The timing constraint \emph{weakens} the ridge, it does not
break it: MRT and $\sinput$ still correlate $-0.64$ to $-0.73$ in the Yasso family, and the
displacement remains a joint one. The claim is that the ridge acquired a slope, not that the two
became independently identified.
\end{warn}

\subsection{Inputs}
\paragraph{The multiplier has a physical referent.} $\sinput \times \bar J$ is an effective litter
flux: $3.19$--$4.34$~tC\,ha$^{-1}$\,yr$^{-1}$ across the six, i.e.\ $18$--$61\%$ above the national
total including understorey ($\sim$2.7). \FIG{SX\_input\_vs\_literature}

\paragraph{Independent corroboration exists.} Liski 2006 --- same country, period, model family and
inventory --- implies a comparable correction. \FIG{F1\_inputs}

\subsection{Kinetics}
\paragraph{Turnover is displaced from published values --- in two of three.} Intrinsic MRT
$24.6$ / $27.1$ / $22.6$~yr (Yasso07/15/20) against published $33.5$ / $30.5$ / $\mathbf{22.1}$.
\FIG{F12\_mrt\_yasso}
\begin{warn}
\textbf{Yasso20's published figure was corrected on 2026-09-02 and the claim changes with it.} The
old $19.0$ came from a vector that paired Yasso20's fractions with \emph{Yasso15's} rates. Corrected,
published Yasso20 is $22.1$ against our $22.6$ --- \textbf{no meaningful displacement}. So the
turnover claim holds for Yasso07 and Yasso15 and \emph{not} for Yasso20, which is consistent with
what the forward experiment then found.
\end{warn}

\paragraph{Not one phenomenon.} Yasso07's whole gap is its \emph{single} climate modifier
($\mathrm{MRT} = \mathrm{MRT}_{\mathrm{ref}}/\xi$ exactly); Yasso15/20 carry three pool-specific
modifiers and cannot use that lever. Do not discuss the family's gaps as one.

\subsection{Why these are ONE claim, not two}
\paragraph{The level pins only the product.} $\mathrm{level} \sim \mathrm{MRT} \times \sinput \times
\bar J$. Faster-with-more and slower-with-less fit identically; the posterior rides that ridge
(Yasso07/15/20 corr $-0.73$/$-0.68$/$-0.64$; the simple models stronger still).
\FIG{F14\_mrt\_ridge} \FIG{S13\_mrt\_ridge\_benchmark}

\begin{warn}
\textbf{The discipline this imposes.} ``Inputs too low'' and ``turnover too slow'' are two
coordinates of \emph{one} displacement. Report it once. Claiming both as independent findings would
double-count, and CLAUDE.md flags this in three places.
\end{warn}

\paragraph{So what \emph{is} the claim?} That the ensemble's best-fitting position on the ridge lies
at higher input and faster turnover than the published parameterisations sit. That is a single,
well-posed statement about location --- and it is testable, because the two ends of the ridge are
\emph{not} equivalent in everything.

\subsection{The independent witness (not on the ridge)}
\paragraph{The models leave an input-shaped signal in their residuals.} They explain $\sim$2\% of SOC
variance; a random forest explains $44$--$65\%$ of what they leave. \FIG{F11\_rf\_heatmap}

\paragraph{And the covariate is the same in all six.} Stand basal area --- a productivity proxy, not
a kinetic one. This does not ride the ridge, so it is genuinely separate evidence that the limitation
is on the input side.

% =====================================================================
\part{Act VII --- Consequences, tested}

\section{If the ensemble sits elsewhere on the ridge, what changes?}
\begin{note}
\textbf{The scientific core, and it has now been run.} A displacement that changed nothing observable
would be bookkeeping. The bet was that it changes the \emph{climate response}. It does --- in
Yasso07, and not in Yasso15 or Yasso20 --- which turns the surprise of Act VI from a curiosity into a
consequence for anyone projecting these soils, and simultaneously bounds how far that consequence
generalises.
\end{note}

\subsection{The hypothesis (as posed, before the run)}
\paragraph{Equal fit, unequal reactivity.} Two calibrations indistinguishable in fit, sitting at
different points on the $\mathrm{MRT} \times \sinput$ ridge, project \emph{different} sensitivity to
warming.

\paragraph{Why it should be so.} Temperature acts on the \emph{rate constants}. A faster-turnover
solution has more of its carbon in pools that feel $\xi$ on a decadal timescale, so the same warming
moves more carbon. The input side responds through productivity instead --- a different pathway, and
one held fixed in these projections.

\paragraph{Already partly visible.} Fit explains only $10$--$18\%$ of forecast spread: $82$--$90\%$
of the projected divergence is among draws the data cannot rank.

\subsection{The proper test: published Yasso vs.\ ours, forward only}
\begin{note}
\textbf{Rejected first: arm A vs arm B.} The $\sinput$ anchor arms look like a two-calibration
experiment but are not a useful one. They share the same model freedom and the same data, and their
optima barely separate --- MRT $22.1$ vs $24.6$~yr, $\sinput$ within $5$--$8\%$. That is a nudge
along the ridge, not the two ends of it.
\end{note}

\paragraph{The real contrast is already published.} The old Yasso parameterisation sits at
\emph{lower input and slower turnover}; our calibration sits at \emph{higher input and faster
turnover}. That is a genuine separation --- Yasso07 MRT $33.5$ with tree-litter-only input, against
$24.6$ with $\sinput = 1.59$: $36\%$ in turnover and $59\%$ in input, three to ten times the arms'.

\paragraph{Forward mode only --- no re-calibration.} Both parameter sets are pushed through the same
engine, same plots, same litter series, same climate, same integrator. Nothing is fitted. This
sidesteps the objection that a re-calibration lets every other parameter readjust, and it is cheap.

\begin{test}
\textbf{Design --- settled, and the published arm is built.} Two axes, one readout.
\smallskip

\begin{center}\small
\begin{tabular}{ll}
\hline
\textbf{Axis 1 --- parameters} & published kinetics, $\sinput$ refitted to the observed window \\
                               & ours (arm B posterior) \\
\textbf{Axis 2 --- climate}    & control (recycled observed, as now) \\
                               & warming ($+2$\,\textdegree C, $+5$\,\textdegree C) \\
\hline
\textbf{Readout} & \multicolumn{1}{l}{\emph{the response}: warming $-$ control, per parameter set} \\
\hline
\end{tabular}
\end{center}
\smallskip
\textbf{$\sinit$ is pinned to our posterior in BOTH arms}, so the two differ in exactly two
coordinates --- turnover and input --- and no initialisation difference is smuggled into a
comparison that is about inputs versus kinetics.
\end{test}

\subsubsection{The published arm agrees with the data. That is the whole point.}
\paragraph{Holding published kinetics fixed, only the input multiplier is refitted.} One parameter,
scored on the three campaign means (balanced set), with $\sinit$ pinned to our posterior in
\emph{both} arms so the contrast is exactly two coordinates on the ridge.
\FIG{doublechecks/published\_arm\_dat.R} \FIG{T\_A1\_published\_arm}

\begin{center}\small
\begin{tabular}{lccccccc}
\hline
& \multicolumn{2}{c}{published arm} & \multicolumn{2}{c}{our arm} & \multicolumn{2}{c}{contrast}
& misfit (pub/ours) \\
& MRT & $\sinput$ & MRT & $\sinput$ & MRT & input & rms log \\
\hline
Yasso07 & 33.47 & 1.249 & 24.57 & 1.585 & $\times0.73$ & $\times1.27$ & 0.020 / 0.029 \\
Yasso15 & 30.54 & 1.335 & 27.11 & 1.538 & $\times0.89$ & $\times1.15$ & 0.022 / 0.034 \\
Yasso20 & 22.13 & 1.724 & 22.56 & 1.659 & $\times1.02$ & $\times0.96$ & 0.029 / 0.035 \\
\hline
\end{tabular}
\end{center}

\paragraph{No compromise was needed.} The published arm matches the observed window \emph{at least
as well as our calibration does} in all three models. So any divergence downstream is extrapolation,
not level mismatch --- exactly the condition the experiment required.

\paragraph{Both arms are on the same ridge.} $\mathrm{MRT} \times \sinput$ is conserved to
$6.9\%$, $2.2\%$ and $2.0\%$. We move \emph{along} each model's ridge, not off it.

\paragraph{The contrast is graded.} Yasso07 is displaced hard, Yasso15 moderately, and
\textbf{Yasso20 hardly at all} --- its published parameterisation already sits where our calibration
puts it, which makes it a within-family control rather than a third replicate.

\begin{note}
\textbf{Two defects in our own code had to be fixed to get here, and both are worth knowing.}
\textbf{(1)} \texttt{assemble\_params} keeps a fixed-rate block, and for Yasso20 those are
\emph{Yasso15's} rates and $w$ parameters --- so a published draw never overrode them. The
``published Yasso20'' MRT was a chimera and F12 plotted it: corrected, the point moves
$19.03 \to 21.85$ and the posterior $25.04 \to 22.13$. \textbf{Yasso20's turnover is therefore not
displaced from its published value at all}, and any claim that it is must be dropped.
\textbf{(2)} Published Yasso20 exhausts the pool budget \emph{exactly} (AWEN fractions plus $p_H$;
median $1.000000$ against Yasso15's $0.99939$), which tripped \texttt{model\_step}'s $0.9999$
diagonal-dominance guard and sent every draw to an Euler fallback that diverges --- a zero-input test
created carbon from nothing. The guard, not the parameters, was wrong: the exact path is stable and
fast at the boundary. Threshold raised to $1.0$, so it fires only when fractions genuinely exceed the
budget. \textbf{Our own posteriors are bit-identical on the new binary}, so no calibrated result
changes.
\end{note}

\subsubsection{Pre-registration, and what happened}
\paragraph{Registered before looking.} Ours should be \emph{more} reactive, with the effect largest
in Yasso07 (single $\xi$ ties turnover and climate sensitivity to one parameter; corr $-0.58$) and
smallest in Yasso15/20 (three pool-specific modifiers decouple them; $-0.10$, $-0.16$). Falsifier:
responses coinciding everywhere, which would mean the ridge is degenerate in climate response too.

\paragraph{Outcome --- SSP2-4.5 at 2084, tC\,ha$^{-1}$, warmed minus stationary.}
\begin{center}\small
\begin{tabular}{lccc}
\hline
& published & ours & difference (90\%) \\
\hline
Yasso07 & $-5.00$ & $-7.08$ & $\mathbf{-2.08\ [-3.57,\ -0.79]}$ \quad excludes zero \\
Yasso15 & $-5.11$ & $-5.50$ & $-0.38\ [-1.30,\ 0.52]$ \\
Yasso20 & $-6.60$ & $-5.82$ & $+0.73\ [-0.42,\ 1.75]$ \\
\hline
\end{tabular}
\end{center}

\paragraph{The registered direction held, but the claim is narrower than registered.} Only Yasso07
separates. Yasso15 does not either, despite a $10\%$ turnover displacement --- so this is \emph{not}
the general statement that ridge position governs reactivity. It is the specific one: \textbf{the
displacement reaches the projection only where a single climate modifier ties turnover and
sensitivity to the same parameter.}

\paragraph{Measured mechanism.} $\mathrm{d}\log\xi_H$ per $+5$\,\textdegree C is $+0.338$ vs
$+0.340$ across Yasso20's arms --- $0.6\%$ --- and $+0.182$ vs $+0.165$ for Yasso15, while Yasso07's
single $\xi$ goes $+0.308 \to +0.412$ and applies to \emph{every} pool. H and N hold $\sim$60\% of
the carbon; the arms differ mainly in the AWE modifier, and AWE holds $\sim$13\%. A steady-state
decomposition predicts the forward runs to a few percent (ratios $1.35/1.05/0.96$ against measured
$1.41/1.03/0.93$), so this is a property of the generators, not of the projection set-up.

\begin{note}
\textbf{The words for it (Lorenzo).} Yasso15 and Yasso20 prove \textbf{resilient to the input-anchor
choice} --- \emph{not} ``right''. This is not a frontal attack on Yasso; it is a measured statement
about which structure is robust to a calibration assumption the data cannot settle.
\end{note}

\begin{warn}
\textbf{Two things that must travel with any number quoted from this experiment.}
\textbf{(1)} The contrast is \emph{paired within a draw} (warmed and stationary share the parameter
vector), which cancels most parameter uncertainty. The unpaired version --- comparing predicted 2084
stocks --- has the same central values but every interval crosses zero, because absolute stock
carries the full within-arm spread. Say which one is being quoted.
\textbf{(2)} Yasso07's published arm is a \emph{point}: no \texttt{.dat} exists, so its interval
carries only our posterior spread while the others convolve two. Its significance is flattered.
\end{warn}

\begin{warn}
$\sinput$ on the published arm is \emph{refitted} to our litter product, not taken from the litter
estimate Yasso was originally calibrated against. The cleanest available choice, but a choice.
\end{warn}

\subsubsection{What we actually publish: one figure in the main text}
\begin{note}
This is an important \emph{side} story, not the spine. One figure in the manuscript, everything
else in the appendix --- and the main-text one leads with the \textbf{consequence}, not the
mechanism.
\end{note}

\paragraph{F15, four panels (BUILT).} Titles \emph{describe} what is plotted; the reading belongs
in the caption and the text, not in the panel heading.
\paragraph{(a) Modelled SOC, climate held stationary.} All three Yassos, both arms
(solid = ours, tight dash = published), $1985 \to 2084$, campaign means with CIs, rule at 2024.
\paragraph{(b) Modelled SOC, warming to SSP2-4.5.} Same models, \textbf{same axes}, projection
warmed on a linear ramp. (a) rises; (b) turns every model over into decline.
\paragraph{(c) Warming response at 2084, by arm.} Boxplots over 200 draws.
\paragraph{(d) Difference between the arms.} Boxplots with a zero line: only Yasso07 excludes zero.

\begin{note}
\textbf{DECIDED 2026-09-02 (Lorenzo).} \textbf{(i) The input anchor is Lehtonen \&
Heikkinen} --- the $\sinput$ prior centre $1.08$, i.e.\ what the run log calls arm~B. Liski's
$1.27$ (arm~A) is retained as a sensitivity, out of the draft, in the notes.
\textbf{(ii) This side arm concludes about the difference in CLIMATE RESPONSE}, not about the
predicted stock. Both are computed --- the prediction gap has the same central values but every
interval crosses zero, because absolute stock carries the full within-arm posterior spread --- and
the response is the claim because it is the paired, within-draw contrast that the data can actually
speak to. Report both; lead with the response.
\end{note}

\begin{note}
\textbf{The result, and the words for it (Lorenzo, 2026-09-02).} Yasso15 and Yasso20 prove
\textbf{resilient to the input-anchor choice} --- \emph{not} ``right''. Their pool-specific climate
modifiers leave the carbon-holding pools (H and N, $\sim$60\% of the stock) responding to warming
almost identically in both arms, while the arms differ mainly in the AWE modifier, and AWE holds
$\sim$13\%. Yasso07 has \emph{one} modifier for every pool, so displacing it displaces the
sensitivity of all the carbon at once. This is not a frontal attack on Yasso; it is a measured
statement about which structure is robust to a calibration assumption the data cannot settle.
\end{note}

\paragraph{Measured mechanism (steady-state decomposition predicts the forward runs).}
$\mathrm{d}\log\xi_H/{+}5$\,\textdegree C is $+0.338$ vs $+0.340$ for Yasso20's two arms --- a
$0.6\%$ difference --- and $+0.182$ vs $+0.165$ for Yasso15. Yasso07's single $\xi$ goes
$+0.308 \to +0.412$, a $34\%$ change applied to every pool. Predicted response ratios
$1.35 / 1.046 / 0.963$ against measured $1.41 / 1.031 / 0.934$.

\paragraph{\textbf{Corollary worth stating in the paper.}} Bulk MRT is the \emph{wrong} summary
statistic for climate sensitivity in a multi-pool model with pool-specific modifiers: two
parameterisations can differ $10\%$ in MRT and be identical in warming response. The right statistic
is the stock-weighted $\sum_i w_i\,\mathrm{d}\log\xi_i/\mathrm{d}T$.
\paragraph{Units and honesty.} All panels in tC\,ha$^{-1}$: the absolute response carries both the
kinetics and the input difference, which is the consequence a reader cares about. The relative
version --- $\sinput$-invariant, therefore pure kinetics --- goes to the appendix. Yasso07's
published arm has no \texttt{.dat}, so it is a \emph{line} where the others have bands, and must be
drawn to look different rather than silently thinner.

\paragraph{Appendix A1 (BUILT).} (a) response against each draw's intrinsic MRT; (b) the same
contrast in relative units; (c) all three SSP scenarios. Plus \texttt{T\_A1}, how the published arm
was constructed.

\begin{note}
\textbf{The collapse test came out as predicted, and it is now a measured result.} A single common
curve (response $\sim \log$ MRT) pooled over all models explains only $R^2 = 0.449$. Within models:
Yasso07 $\mathbf{0.703}$ --- its single $\xi$ ties turnover and climate sensitivity to the same
parameter --- against Yasso15 $0.209$ and Yasso20 $\mathbf{0.068}$, where three pool-specific
modifiers decouple them. Visually the three clouds do not lie on one curve: Yasso20 at
MRT $\approx 22$ responds $-8.8\%$ while Yasso07 draws at the \emph{same} MRT respond $\approx
-10\%$. \textbf{Bulk MRT does not determine the warming response}, which is the evidence for the
corollary below.
\end{note}

\begin{note}
\textbf{Why the collapse figure matters even in an appendix.} It is what converts ``two
parameterisations differ'' into ``position on the ridge governs reactivity''. If the clouds do not
collapse onto one curve, the effect is model identity rather than ridge position --- a different
claim, and one this design would not support.
\end{note}

\NEW{F15 (main text, 3 panels) + appendix: MRT collapse, all-model trajectories, arm table}

\subsection{Stronger versions, if this one pays off}
\paragraph{A real scenario rather than a temperature offset.} Replace the $+2$/$+5$\,\textdegree C
construct with a proper climate projection, so the answer is a projection and not a sensitivity.

\paragraph{Let productivity respond too.} Warming acts on inputs as well as on decomposition; here
the input side is held fixed, which biases the comparison toward loss in \emph{both} arms. Coupling
it is the honest version and a much larger job.

\paragraph{Break the degeneracy from outside.} Radiocarbon constrains turnover independently of
input magnitude. The degeneracy is a property of the observation design, not of nature. Worth
stating even if we cannot do it here.

\begin{warn}
\textbf{The projections currently contain no climate change} --- the last twenty observed years are
recycled. Everything above is therefore about \emph{sensitivity}, computed separately, and the
forecast figures are disequilibrium relaxation under stationary climate. The paper must say so.
\end{warn}

% =====================================================================
\part{Act VIII --- Thread C}

\section{What the data can and cannot constrain}
\paragraph{Convergence is not identifiability.} R-hat $\le 1.008$ everywhere, yet the transfer
fractions are prior-pinned, ridge-correlated and carry $<1$ nat. A clean sampler on a flat ridge.
\FIG{F8b\_fraction\_ridges} \FIG{T2\_fraction\_ridges}

\paragraph{The data pin a level, and little else.} Two observations per plot; no cross-plot skill.
\FIG{F6\_obs\_vs\_pred\_holdout}

\paragraph{Report it, do not fix it.} Non-identifiability here is a property of the observation
design and belongs in the results.

% =====================================================================
\part{Act IX --- Landing}

\section{The policy deliverable}
\paragraph{An estimate plus its irreducible structural uncertainty.} Not a single number, and not a
selected winner.

\paragraph{Aggregate, not plot-level.} These models capture the national trajectory; they do not
rank plots. Say it early --- otherwise ``captures the dynamic'' and ``$R^2 \approx 0$'' read as a
contradiction.

\paragraph{Maps as the landing figure.} \FIG{NextGenC ensemble maps}

% =====================================================================
\part{Act X --- Outlook}

\section{Where the recoverable signal hides}
\paragraph{Historical inputs.} The deepest uncertainty; $\sinit$ stands in for a century.
\paragraph{Understorey and belowground.} The named gap in the litter product.
\paragraph{Stratified inputs.} The understorey fraction rises strongly south to north; one national
multiplier cannot carry it. \textbf{Explicitly NOT part of this paper} (decided 2026-09-02) --- the
stratified-calibration strand is dropped, kept behind a switch in the manuscript source in case it
is picked up later.
\paragraph{Radiocarbon.} The clean external constraint on turnover.

% =====================================================================
\part{Appendix}
\section{Figure inventory}
\paragraph{Current.} \textbf{All 31 builders re-run on arm B} (\texttt{20260831\_1624*}) on
2026-09-02, zero failures --- F1--F14, S2--S14, R3/R4, T1, T2, plus the new \textbf{F15}
(four panels) and \textbf{A1}/\textbf{T\_A1} (the forward-experiment appendix).
\paragraph{Fixed in the same pass.} T1's $\bar J$ constant ($2.472 \to 2.511$); the multimodel
metrics table, which had been reporting all-data rather than calibration-only skill and now agrees
with T1 for the first time; F5's campaign mapping and rate denominator; and F12's published Yasso20.
\paragraph{Still to build.} F3b (the three trajectory-shape groups of Act IV).
\paragraph{Not started.} Citations, a \texttt{.bib}, figure placement in the manuscript, and the
title/abstract/coauthor block --- deliberately held until this storyline is settled.

\end{document}
